% -------------------------------------------------------------------------
% WBS ID: RES-NUM-FRM
% Deliverable: Numerical Formulation and Discrete Conservation Framework
% Status: Draft complete
% Evidence status: Regional and local finite-volume research specifications complete at family level
% Review status: Not reviewed
% Last updated: 2026-07-18
% Dependencies: RES-MOD-EQS, RES-MOD-SRC
% Downstream interfaces: RES-NUM-SPEC, RES-NUM-IBC, RES-NUM-STB, Software finite-volume implementation
% -------------------------------------------------------------------------

\section{RES-NUM-FRM --- Numerical Formulation and Discrete Conservation Framework}
\label{sec:res-num-frm}

\subsection{Purpose and Scope}
\label{sec:res-num-frm-purpose}

This Deliverable records the breadth-first numerical-method decision
for the regional tsunami solver. It selects the numerical family used
for the two-dimensional nonlinear shallow-water model and records the
candidate families that remain useful as references, deferred ideas, or
screened alternatives.

The original family-level decision selected finite volume for the
regional model. The 2026-07-07 regional gate extends that decision to a
research-specification-level numerical model for the regional 2D NLSWE
solver. The local three-dimensional CFD family is also finite volume,
but it is documented separately from the regional NLSWE discretisation
because it solves a collocated, polyhedral, pressure-corrected
URANS--VOF system.

\subsection{Research Questions}
\label{sec:res-num-frm-research-questions}

The current phase answers the following questions:

\begin{enumerate}
  \item Which numerical family is selected for the production regional
    nonlinear shallow-water solver?
  \item Which candidate families are retained as references or
    deferred concepts?
  \item Which lower-level choices remain unresolved after the
    family-level decision?
\end{enumerate}

\subsection{Terminology and Definitions}
\label{sec:res-num-frm-definitions}

Four decision statuses are used:

\begin{description}
  \item[selected] Adopted for the current project branch.
  \item[provisional] Adopted as a working direction pending lower-level
    comparison or verification.
  \item[excluded] Removed from the active branch for project-specific
    reasons.
  \item[unresolved] Not yet selected, excluded, or formally deferred.
\end{description}

\subsection{Literature Evidence}
\label{sec:res-num-frm-literature-evidence}

% TODO: Record evidence from peer-reviewed papers, authoritative datasets,
% standards, technical reports and primary documentation.
%
% Do not create a detached literature-summary list. Explain how each source
% supports, contradicts or limits a project decision.

\subsubsection{Finite-Difference Method}
\label{sec:res-num-frm-evidence-fdm}

\paragraph{Methodological basis.}
\textcite{casulli1992semiimplicit} developed a semi-implicit finite
difference architecture for shallow-water flow in which the free-surface
unknowns are reduced to symmetric positive-definite linear systems. The
paper is retained for selective implicit treatment, conjugate-gradient
and multigrid compatibility, stable wetting and drying, and the
distinction between stability-limited and accuracy-limited time steps.

\paragraph{Contextual validation.}
\textcite{wei2013japan} used MOST to model the 2011 Japan tsunami and
demonstrated a finite-difference family in a near-field tsunami
forecast setting. This validates a production tsunami use case for the
family, not the exact Casulli--Cheng semi-implicit formulation.

\subsubsection{Finite-Volume Method}
\label{sec:res-num-frm-evidence-fvm}

\paragraph{Methodological basis.}
\textcite{kurganov2007wellbalanced} presented a second-order
central-upwind method for the Saint-Venant system that simultaneously
targets well-balancing and positivity preservation. The method is
retained as primary evidence for the regional need to preserve lake-at-
rest balance, non-negative depth, discontinuities, and dry or nearly dry
states.

\paragraph{Contextual validation.}
\textcite{dutykh2011volna} documented the VOLNA tsunami code using
finite volumes on unstructured triangular meshes for tsunami generation,
propagation, run-up and inundation. The paper is retained as contextual
evidence that a finite-volume shallow-water solver can support complex
coastal domains, wetting and drying, and tsunami-scale validation cases.

\paragraph{Local 3D finite-volume and pressure-correction evidence.}
\textcite{hirt1981vof} established the VOF representation of moving
free boundaries. \textcite{vukcevic2018sharp} and
\textcite{roenby2016computational} support sharp and geometric VOF
treatments on finite-volume meshes. \textcite{issa1986operator}
supports pressure-correction by operator splitting for incompressible
flow. Source roles: `3D-VOF` and `3D-FVM`. Project conclusion: the
local solver shall use conservative finite-volume transport for mass,
momentum and phase fraction, with pressure correction enforcing
incompressibility. Limitation: these sources establish the numerical
family and pressure-correction principle, not the final linear solver,
preconditioner or mesh-convergence tolerance.

\subsubsection{Lattice/Discrete Boltzmann Method}
\label{sec:res-num-frm-evidence-dbm}

\paragraph{Methodological basis.}
\textcite{meng2018discreteboltzmann} derived discrete Boltzmann models
for shallow-water equations using polynomial equilibria. The paper is
retained for collision-level conservation, local stream--collision
structure, GPU-relevant locality, and the conditional ability to improve
supercritical capability through higher-order expansion and quadrature.

\paragraph{Contextual validation gap.}
No tsunami-context validation paper for the lattice or discrete
Boltzmann branch has yet passed the project evidence gate. The reviewed
paper does not establish genuinely dry wetting and drying, fully
discrete well-balancing, positivity of depth in the project setting, or
adequate treatment of irregular coastlines and bathymetry.

\subsubsection{Spectral-Method Screening}
\label{sec:res-num-frm-evidence-spectral-screening}

No detailed spectral-method literature-evidence subsection is created
for the current comparison. Spectral and pseudospectral methods are
screened at the project-fit level owing to non-periodic irregular
coastal boundaries, wetting and drying, bores and hydraulic
discontinuities, complex bathymetry, local refinement requirements, and
the loss of the principal global-spectral advantages once extensive
boundary treatment, filtering, and domain decomposition are required.

\subsection{Governing Relations and Quantitative Definitions}
\label{sec:res-num-frm-governing-relations}

% TODO: Record relevant equations, dimensionless groups, empirical models,
% extraction rules or quantitative definitions.
%
% Use align environments for displayed equations.
% Every displayed equation must have a unique \label{eq:...}.
% Do not place punctuation after displayed equations.

\subsection{Discrete Formulation}
\label{sec:res-num-frm-discrete-formulation}

The regional production solver is selected as a cell-centred
finite-volume method on unstructured triangular or polygonal control
volumes. This selection concerns the spatial conservation framework, not
the full lower-level numerical method.

Spatial discretisation and source-term integration are separate
decisions. The regional nonlinear shallow-water equations already
include bed slope, bottom friction, Coriolis forcing, earthquake or
moving-bottom forcing where activated, and wetting--drying
requirements. The provisional source-term treatment is:

\begin{description}
  \item[bed slope] coupled well-balanced hydrostatic reconstruction;
  \item[bottom friction] semi-implicit update;
  \item[Coriolis] exact or structure-preserving rotational update.
\end{description}

The detailed derivations and notation for friction and Coriolis remain
owned by the mathematical-model Deliverables. This numerical
Deliverable records how those terms are integrated within the selected
finite-volume branch.

\subsubsection{Local 3D Collocated Polyhedral Finite Volume}
\label{sec:res-num-frm-local-polyhedral-finite-volume}

The local URANS--VOF solver uses cell-centred collocated fields on
polyhedral control volumes. For a local cell $K_i$ with volume $V_i$,
faces $f$, outward area vectors $\mathbf{S}_f$ and neighbouring cells
$\mathcal{N}(i)$, a transported scalar or vector component $\psi$ is
stored as:

\begin{align}
  \psi_i(t)
  &=
  \frac{1}{V_i}
  \int_{K_i}
  \psi(\mathbf{x},t)
  \,\dd V
  \label{eq:res-num-frm-local-cell-average}
\end{align}

The face-normal volume flux is:

\begin{align}
  \phi_f
  &=
  \mathbf{u}_f
  \cdot
  \mathbf{S}_f
  \label{eq:res-num-frm-local-face-volume-flux}
\end{align}

The conservative semi-discrete transport form is:

\begin{align}
  \frac{\dd}{\dd t}
  \left(
    V_i\psi_i
  \right)
  &=
  -
  \sum_{f\in\partial K_i}
  \phi_f\psi_f
  +
  \sum_{f\in\partial K_i}
  \Gamma_f
  \nabla\psi_f
  \cdot
  \mathbf{S}_f
  +
  V_i S_{\psi,i}
  \label{eq:res-num-frm-local-conservative-transport}
\end{align}

After temporal and spatial linearisation each transported field is
assembled into an algebraic cell equation:

\begin{align}
  a_{P,i}\psi_i
  +
  \sum_{j\in\mathcal{N}(i)}
  a_{N,ij}\psi_j
  &=
  b_i
  \label{eq:res-num-frm-local-algebraic-cell-equation}
\end{align}

Source terms may be linearised as:

\begin{align}
  S_{\psi,i}
  &=
  S_{C,i}
  +
  S_{P,i}\psi_i
  \label{eq:res-num-frm-local-source-linearisation}
\end{align}

with coefficient updates:

\begin{align}
  a_{P,i}^{\mathrm{new}}
  &=
  a_{P,i}^{\mathrm{base}}
  -
  V_i S_{P,i}
  \label{eq:res-num-frm-local-source-ap-update}
  \\
  b_i^{\mathrm{new}}
  &=
  b_i^{\mathrm{base}}
  +
  V_i S_{C,i}
  \label{eq:res-num-frm-local-source-b-update}
\end{align}

Pressure correction enforces the discrete incompressibility constraint:

\begin{align}
  \sum_{f\in\partial K_i}
  \phi_f^{\,n+1}
  &=
  0
  \label{eq:res-num-frm-local-discrete-continuity}
\end{align}

The provisional pressure-correction form is:

\begin{align}
  \mathbf{u}^{\,n+1}
  &=
  \mathbf{u}^{*}
  -
  \mathbf{D}
  \nabla p'
  \label{eq:res-num-frm-local-velocity-correction}
  \\
  p^{\,n+1}
  &=
  p^{*}
  +
  p'
  \label{eq:res-num-frm-local-pressure-correction}
\end{align}

where $\mathbf{u}^{*}$ is the momentum-predictor velocity,
$p'$ is the pressure correction and $\mathbf{D}$ is the discrete
velocity-correction operator implied by the momentum matrix. Rhie--Chow
or an equivalent pressure-consistent face-flux interpolation shall be
used on collocated meshes if required to prevent pressure--velocity
decoupling.

\subsection{Conservation Properties}
\label{sec:res-num-frm-conservation-properties}

% TODO: Complete this RES-NUM Domain-specific subsection.

\subsection{Consistency and Formal Accuracy}
\label{sec:res-num-frm-consistency-formal-accuracy}

% TODO: Complete this RES-NUM Domain-specific subsection.

\subsection{Stability Requirements}
\label{sec:res-num-frm-stability-requirements}

% TODO: Complete this RES-NUM Domain-specific subsection.

\subsection{Mesh and Time-Step Requirements}
\label{sec:res-num-frm-mesh-time-step-requirements}

% TODO: Complete this RES-NUM Domain-specific subsection.

\subsection{Numerical Failure Modes}
\label{sec:res-num-frm-numerical-failure-modes}

% TODO: Complete this RES-NUM Domain-specific subsection.

\subsection{Verification Requirements}
\label{sec:res-num-frm-verification-requirements}

% TODO: Complete this RES-NUM Domain-specific subsection.

\subsection{Candidate Approaches}
\label{sec:res-num-frm-candidate-approaches}

The candidate regional numerical families considered in this phase are:

\begin{itemize}
  \item cell-centred finite volume;
  \item finite difference;
  \item lattice or discrete Boltzmann;
  \item spectral or pseudospectral methods.
\end{itemize}

\subsection{Critical Comparison}
\label{sec:res-num-frm-critical-comparison}

% TODO: Compare candidates using physically and project-relevant criteria.
% Do not select an approach solely on popularity or implementation ease.

\subsubsection{Comparative Assessment}
\label{sec:res-num-frm-comparative-assessment}

The comparison is project-specific rather than a universal ranking of
numerical analysis families.

\begin{description}
  \item[Conservation] Finite volume is strongest for direct
    conservation of water depth and horizontal momentum through face
    fluxes. Finite difference remains useful but requires careful
    conservative formulation. Lattice/discrete Boltzmann has strong
    collision-level conservation, while spectral conservation depends on
    the formulation and filtering.
  \item[Well-balancing] The selected finite-volume branch has direct
    support through hydrostatic reconstruction and well-balanced shallow-
    water fluxes. Finite-difference and Boltzmann branches retain useful
    ideas but require further project-specific proof.
  \item[Positivity] Finite volume provides the clearest route to
    non-negative water depth in wetting--drying regions. The reviewed
    Boltzmann paper does not establish the required dry-front
    robustness.
  \item[Wetting and drying] VOLNA and the central-upwind evidence
    support finite volume for inundation. Finite difference remains
    credible as a reference family through MOST and semi-implicit
    shallow-water literature.
  \item[Bore and shock treatment] Finite volume is selected for weak
    solutions, bores, hydraulic discontinuities, and shoreline Riemann
    treatments.
  \item[Coastline and mesh flexibility] Unstructured finite volumes
    match irregular coastlines, harbour geometry, local mesh grading,
    and face-based boundary patches.
  \item[Phase and amplitude accuracy] The exact regional flux,
    reconstruction, timestep and optional dispersive extension remain
    unresolved and shall be tested against propagation and inundation
    benchmarks.
  \item[Time-step and linear-system structure] Finite difference is
    retained for selective implicitness, reduced free-surface systems,
    symmetric positive-definite solves, conjugate-gradient compatibility
    and multigrid compatibility.
  \item[Parallel and GPU structure] Finite volume and lattice/discrete
    Boltzmann both expose local operations. Lattice/discrete Boltzmann
    remains deferred until the tsunami-validation and wet/dry evidence
    gaps are closed.
  \item[Coupling compatibility] Finite volume produces conservative
    face fluxes suitable for the two-dimensional--three-dimensional
    interface and shares cell, face, flux and boundary-patch abstractions
    with the OpenFOAM-inspired local solver architecture.
\end{description}

\subsection{Assumptions and Applicability}
\label{sec:res-num-frm-assumptions}

% TODO: State assumptions and the conditions under which they are valid.

\subsection{Limitations and Failure Conditions}
\label{sec:res-num-frm-limitations}

% TODO: State where the evidence, model or selected approach ceases to be
% reliable or applicable.

The finite-volume selection does not imply that finite difference,
lattice/discrete Boltzmann, or spectral methods are inferior in general.
The screening is limited to the current regional tsunami solver branch.

Finite difference is not selected for production regional execution but
is retained as a numerical reference and as a source of semi-implicit
architectural ideas. Lattice/discrete Boltzmann is removed from the
active regional shortlist until tsunami validation, positivity,
well-balancing, dry wetting and drying, and complex-bathymetry evidence
are available. Spectral methods are screened out of the active regional
comparison owing to project-fit constraints rather than mathematical
invalidity.

\subsection{Project-Specific Conclusions}
\label{sec:res-num-frm-conclusions}

% TODO: Convert the evidence into conclusions specific to the Studentship
% project. Do not merely restate literature findings.

The regional nonlinear shallow-water solver shall use a cell-centred
finite-volume method. The primary rationale is physical and numerical:
direct conservation, weak-solution compatibility, well-balanced
bathymetry, positivity-preserving wetting and drying, unstructured
coastline representation, local mesh grading, and conservative
two-dimensional--three-dimensional interface fluxes.

OpenFOAM compatibility is a secondary architectural benefit. The
regional solver is not required to be implemented inside OpenFOAM.

The local three-dimensional solver shall also use finite volume, but
with collocated polyhedral control volumes, VOF phase transport and
pressure correction for incompressibility. OpenFOAM is retained as the
principal architectural reference for this local branch.

\subsection{Selected Approach and Rationale}
\label{sec:res-num-frm-selected-approach}

% TODO: Record the selected approach only after sufficient evidence exists.
% State the alternatives rejected and the reasons for rejection.

\subsubsection{Regional Solver Decision}
\label{sec:res-num-frm-regional-solver-decision}

The selected regional numerical family is:

\begin{quote}
  cell-centred finite-volume method
\end{quote}

The selected regional solver composition is:

\begin{itemize}
  \item unstructured triangular or polygonal control volumes;
  \item second-order MUSCL reconstruction;
  \item weighted least-squares gradients;
  \item Barth--Jespersen limiting;
  \item well-balanced hydrostatic reconstruction;
  \item positivity-preserving wet/dry treatment;
  \item HLLC primary flux with HLL/HLLE fallback near dry or unsafe
    states;
  \item SSPRK(3,3) transport;
  \item Strang-split semi-implicit Manning friction;
  \item characteristic or radiation open boundaries;
  \item sponge damping for artificial ocean cuts where required;
  \item mesh refinement driven by bathymetry, shoreline, source,
    structures, wet/dry fronts and validation observables.
\end{itemize}

Dry-depth tolerances, exact wave-speed estimates, final mesh sizes and
domain-specific sponge dimensions remain verification or implementation
parameters, not open research-family decisions.

\subsubsection{Retained FDM Concepts}
\label{sec:res-num-frm-retained-fdm-concepts}

Finite difference is retained for:

\begin{itemize}
  \item selective implicit treatment;
  \item reduced free-surface systems;
  \item symmetric positive-definite linear systems;
  \item conjugate-gradient and multigrid compatibility;
  \item stable wet/dry face handling;
  \item the distinction between stability-limited and accuracy-limited
    timesteps.
\end{itemize}

\subsubsection{Excluded and Deferred Alternatives}
\label{sec:res-num-frm-excluded-deferred-alternatives}

Lattice/discrete Boltzmann is removed from the active regional
shortlist. It remains a deferred research candidate if tsunami-context
validation, robust dry wetting and drying, fully discrete
well-balancing, and irregular-coastline treatment are demonstrated.

Spectral and pseudospectral methods are screened out of the active
regional solver comparison for this project branch.

\subsubsection{Local 3D Solver Decision}
\label{sec:res-num-frm-local-solver-decision}

The selected local numerical family is:

\begin{quote}
  cell-centred collocated polyhedral finite volume with bounded VOF
  phase transport and pressure correction for incompressible
  two-phase URANS
\end{quote}

The selected baseline turbulence closure for this finite-volume branch
is $k$--$\omega$ SST. LES remains a higher-fidelity comparison mode,
and particle, finite-element and staggered-grid methods remain
references or specialist alternatives rather than the production local
hydrodynamic family.

\subsection{Unresolved Decisions}
\label{sec:res-num-frm-unresolved-decisions}

% TODO: Record unresolved questions, required evidence and decision owners.

The following regional implementation and validation parameters remain
open:

\begin{itemize}
  \item dry-depth tolerance;
  \item exact HLL/HLLC wave-speed safeguards;
  \item final mesh refinement levels;
  \item final CFL value after positivity and wet/dry verification;
  \item optional dispersive extension and its activation gate;
  \item domain-specific open-boundary and sponge placement;
  \item coupling residuals and convergence criteria for later
    multidimensional coupling.
  \item exact local VOF advection and compression scheme;
  \item exact pressure-correction sequence and pressure-consistent
    face-flux interpolation;
  \item local wall-function and rough-wall treatment;
  \item local linear solvers, preconditioners and residual tolerances.
\end{itemize}

\subsection{Requirements and Handoffs}
\label{sec:res-num-frm-requirements}

% TODO: State requirements passed to other Research Domains, Software,
% verification activities or later execution work.

The selected regional solver family hands off the following
requirements:

\begin{itemize}
  \item `RES-NUM-MSH` shall define unstructured mesh, local grading and
    adaptive-mesh strategy;
  \item `RES-NUM-FLX` shall compare HLL, HLLC, central-upwind and
    related shoreline-aware fluxes;
  \item `RES-NUM-SRC` shall formalise bed-slope, friction and Coriolis
    integration;
  \item `RES-NUM-TIM` shall select the SSP-RK order and any
    semi-implicit split;
  \item `RES-NUM-WET` shall select wet/dry thresholds and positivity
    controls;
  \item `RES-MOD-MUL` and later Software work shall use conservative
    face-flux outputs for the two-dimensional--three-dimensional
    interface.
  \item `RES-NUM-SPEC` shall define the local URANS--VOF solve order,
    pressure-correction loop, residual checks and timestep rejection
    rules.
\end{itemize}

\subsection{Deliverable Completion Record}
\label{sec:res-num-frm-completion-record}

% TODO: Record whether the Deliverable is:
%
% - Not started
% - In progress
% - Draft complete
% - Reviewed
% - Accepted
%
% Acceptance requires:
%
% - the research questions to be answered;
% - claims to be supported by appropriate evidence;
% - assumptions and limitations to be explicit;
% - the project-specific conclusion to be stated;
% - downstream requirements to be traceable;
% - unresolved decisions to be closed or formally deferred.
